\subsection{Support Vector Machines and Margin Methods}

Margin methods prefer classifiers that separate classes with a buffer, not merely with zero training error.

\begin{definition} \label{def:sl-functional-margin} \hlt{Margin}\\
For labels $y_i\in\{-1,1\}$ and score $g(x)=w^\top x+b$, the signed margin of $(x_i,y_i)$ is $y_ig(x_i)$.  Large positive margins indicate confident correct classification.
\end{definition}

\begin{definition} \label{def:sl-soft-margin-svm} \hlt{Soft-Margin Support Vector Machine}\\
The linear soft-margin SVM solves
\[
\min_{w,b,\xi}
\left\{\frac12\norm{w}_2^2+C\sum_{i=1}^n\xi_i\right\}
\]
subject to $y_i(w^\top x_i+b)\geq1-\xi_i$ and $\xi_i\geq0$.  Equivalently, it minimizes an $\ell^2$ penalty plus hinge loss.
\end{definition}

\begin{remark} \label{rem:sl-support-vectors} \hlt{Support Vectors}\\
Only examples on or inside the margin determine the fitted separating hyperplane.  In kernel form, these examples are the nonzero coefficients in the representer expansion.
\end{remark}

\begin{remark} \label{rem:sl-margin-generalization} \hlt{Margins and Generalization}\\
Margin bounds refine raw VC counting by using the geometry of the fitted classifier.  The effective complexity depends on both norm control and the achieved margin.
\end{remark}
